% Experiment 15: Implementation and Security Analysis of Knapsack Cryptosystem
\section{Theory}
The Merkle-Hellman Knapsack Cryptosystem is an early public-key scheme based
on the NP-complete subset-sum problem. The private key is a \emph{superincreasing}
sequence $(b_1, b_2, \dots, b_n)$ where each element exceeds the sum of all
preceding elements. A modulus $m > \sum b_i$ and coprime multiplier $w$ are
chosen, and the public key is $a_i = (w \cdot b_i) \bmod m$. Encryption of
an $n$-bit plaintext block $(x_1, \dots, x_n)$ is a subset sum:
\[
C = \sum_{i=1}^{n} x_i a_i
\]
Decryption transforms the ciphertext into the superincreasing domain using
$C' = (C \cdot w^{-1}) \bmod m$, after which the original bits are recovered
efficiently with a greedy algorithm.

\section{Algorithm}
\begin{enumerate}
  \item Generate a private superincreasing sequence $(b_1, \dots, b_n)$.
  \item Choose modulus $m$ greater than $\sum b_i$.
  \item Choose multiplier $w$ such that $\gcd(w, m) = 1$.
  \item Compute the public key $a_i = (w \cdot b_i) \bmod m$ for each $i$.
  \item To encrypt an $n$-bit plaintext block: compute $C = \sum x_i a_i$ using the public key.
  \item To decrypt: compute $w^{-1} \bmod m$ and $C' = (C \cdot w^{-1}) \bmod m$.
  \item Solve the superincreasing subset-sum for $C'$ using the private sequence $(b_1,\dots,b_n)$ with a greedy algorithm (from largest to smallest).
  \item Recover the original plaintext bits and verify correctness.
\end{enumerate}

\section{Code}
\begin{minted}{cpp}
#include<bits/stdc++.h>
using namespace std;
class knapsack {
  private:
    int n;
    vector<long long> priv, pub;
    long long m, w;
    long long gcd(long long a, long long b) {
      return b == 0 ? a : gcd(b, a % b);
    }
    long long modInverse(long long a, long long mod) {
      long long m0 = mod, x0 = 0, x1 = 1;
      a = ((a % mod) + mod) % mod;
      while (a > 1) {
        long long qt = a / mod;
        long long t = mod;
        mod = a % mod;
        a = t;
        t = x0;
        x0 = x1 - qt * x0;
        x1 = t;
      }
      if (x1 < 0) x1 += m0;
      return x1;
    }
  public:
    knapsack(int bits);
    long long encrypt(string bits);
    string decrypt(long long cipher);
};
knapsack::knapsack(int bits) {
  n = bits;
  priv.resize(n);
  long long sum = 0;
  for (int i = 0; i < n; i++) {
    priv[i] = sum + rand() % 5 + 1;
    sum += priv[i];
  }
  cout << "Private (superincreasing) sequence: ";
  for (auto v : priv) cout << v << " ";
  cout << endl;

  m = sum + rand() % 10 + 1;
  do {
    w = 2 + rand() % (int)(m - 2);
  } while (gcd(w, m) != 1);
  cout << "Modulus m = " << m << ", Multiplier w = " << w << endl;

  pub.resize(n);
  for (int i = 0; i < n; i++) pub[i] = (priv[i] * w) % m;
  cout << "Public key sequence: ";
  for (auto v : pub) cout << v << " ";
  cout << endl;
}
long long knapsack::encrypt(string bits) {
  long long cipher = 0;
  for (int i = 0; i < n; i++)
    if (bits[i] == '1') cipher += pub[i];
  return cipher;
}
string knapsack::decrypt(long long cipher) {
  long long wInv = modInverse(w, m);
  long long cPrime = (cipher * wInv) % m;
  cout << "w^-1 mod m = " << wInv << endl;
  cout << "C' = C * w^-1 mod m = " << cPrime << endl;

  string decrypted(n, '0');
  for (int i = n - 1; i >= 0; i--) {
    if (priv[i] <= cPrime) {
      decrypted[i] = '1';
      cPrime -= priv[i];
    }
  }
  return decrypted;
}
int main() {
  cout << "Enter number of bits in the block (e.g., 8): ";
  int n;
  cin >> n;

  knapsack cipher(n);

  cout << "\nEnter a " << n << "-bit binary plaintext block: ";
  string bits;
  cin >> bits;

  long long cipherText = cipher.encrypt(bits);
  cout << "\n--- Encryption ---" << endl;
  cout << "Ciphertext (sum using public key) = " << cipherText << endl;

  cout << "\n--- Decryption ---" << endl;
  string decrypted = cipher.decrypt(cipherText);
  cout << "Decrypted bits: " << decrypted << endl;
  cout << "Original bits : " << bits << endl;

  cout << "\nVerification: "
       << (decrypted == bits ? "SUCCESS - Plaintext recovered correctly"
                              : "FAILURE - Mismatch")
       << endl;
  return 0;
}
\end{minted}

\section{Result and Discussion - I/O}
\begin{verbatim}
Enter number of bits in the block (e.g., 8): 8
Private (superincreasing) sequence: 4 6 13 24 51 99 199 399
Modulus m = 805, Multiplier w = 219
Public key sequence: 71 509 432 426 704 751 111 441

Enter a 8-bit binary plaintext block: 10101100

--- Encryption ---
Ciphertext (sum using public key) = 1958

--- Decryption ---
w^-1 mod m = 669
C' = C * w^-1 mod m = 167
Decrypted bits: 10101100
Original bits : 10101100
\end{verbatim}

\textbf{Observation:} The public key does not appear superincreasing, which
disguises the private structure; decryption uses $w^{-1}$ to restore the
easy domain where a greedy algorithm recovers the bits.

\section{Conclusion}
The Merkle-Hellman Knapsack cryptosystem was implemented, successfully
encrypting and decrypting an 8-bit block. This scheme is now considered
broken (Shamir showed the transformation leaks enough structure to recover
the private key) and is retained for educational purposes.
